\documentclass[11pt,a4paper]{article}

\usepackage[margin=1in]{geometry}
\usepackage{amsmath,amssymb,amsthm}
\usepackage{booktabs}
\usepackage{hyperref}
\usepackage{xcolor}
\usepackage{enumitem}
\usepackage{microtype}
\usepackage{parskip}
\usepackage{titlesec}
\usepackage{array}
\usepackage{longtable}

\hypersetup{colorlinks=true, linkcolor=blue, urlcolor=blue, citecolor=blue}

\title{\textbf{DSM2: Distilled Sequence Model v2}\\
\large Technical Report of the Training Pipeline}
\author{GleghornLab}
\date{\today}

\begin{document}
\maketitle
\tableofcontents
\newpage

% ============================================================
\section{Introduction}
% ============================================================

DSM2 (Distilled Sequence Model v2) is a self-supervised protein language model (pLM) trained via
knowledge distillation from a frozen, large-scale teacher pLM.  The student is an E1-architecture
transformer that learns to (i) reconstruct masked amino-acid tokens (masked language modelling,
MLM), (ii) match the intermediate hidden-state representations of the teacher on a per-token basis
(JEPA loss), and (iii) preserve the intra-batch relational geometry of the teacher's representations
at the sequence level (contrastive loss).

The training pipeline has three sequential loss phases, an exponential-moving-average (EMA) teacher
that eventually replaces the external teacher, and supports both standard AdamW and the MuonClip
optimizer (Muon + QK-Clip).  Distributed data-parallel (DDP) training is supported out of the box.

The default teacher is \texttt{Synthyra/DPLM-3B}; the default student has a hidden dimension of
256, trained for 100\,000 optimiser steps on \texttt{Synthyra/uniref50}.

% ============================================================
\section{Model Architecture}
% ============================================================

\subsection{Teacher Model}

The teacher is a frozen, pre-trained protein language model loaded in \texttt{bfloat16} and moved to
the target device.  All teacher parameters have \texttt{requires\_grad = False}; the teacher is kept
in \texttt{eval()} mode throughout training.  The teacher is selected by inspecting the model path
string:

\begin{center}
\begin{tabular}{ll}
\toprule
\textbf{Path substring} & \textbf{Class loaded} \\
\midrule
\texttt{dplm2} & \texttt{DPLM2Model} \\
\texttt{dplm}  & \texttt{DPLMModel} \\
\texttt{esm2}, \texttt{fastesm} & \texttt{FastEsmModel} \\
\texttt{e1}    & \texttt{E1ForMaskedLM} \\
(default)      & \texttt{ESMplusplusModel} \\
\bottomrule
\end{tabular}
\end{center}

The teacher exposes \texttt{config.num\_hidden\_layers} (number of transformer layers) and
\texttt{config.hidden\_size} ($d_T$, teacher hidden dimension), both of which are used to configure
the student.

\subsection{Student Model (DSM2 / E1 Backbone)}

The student model is an instance of \texttt{DSM2}, a subclass of \texttt{E1ForMaskedLM} augmented
with distillation infrastructure.  The E1 backbone is a decoder-style transformer with pre- and
post-attention RMSNorm, grouped-query attention (GQA), and a SwiGLU (or standard) MLP.

\subsubsection{Configuration (\texttt{DSM2Config})}

\texttt{DSM2Config} extends \texttt{E1Config} with two additional fields:

\begin{itemize}
  \item \texttt{teacher\_hidden\_size} ($d_T$): hidden dimension of the teacher (default 768).
  \item \texttt{use\_teacher\_projections}: whether to attach a linear projection per layer that
        maps student representations into teacher space (default \texttt{True}; disabled after the
        EMA teacher takes over).
\end{itemize}

The student is initialised from scratch when \texttt{pretrained\_weights} is \texttt{None}; the
following derived quantities are computed from the CLI arguments:

\begin{align*}
  d_s &= \texttt{student\_hidden\_size} \quad (\text{must satisfy } d_s \bmod 64 = 0) \\
  d_{\mathrm{ff}} &= \lfloor d_s \times \texttt{student\_expansion\_ratio} \rfloor \\
  H   &= d_s / 64 \quad (\text{number of attention heads, also number of KV heads})\\
  N   &= \texttt{teacher.config.num\_hidden\_layers} \quad (\text{shared with teacher})
\end{align*}

\subsubsection{E1 Transformer Backbone}

Each of the $N$ transformer layers consists of:

\begin{enumerate}
  \item \textbf{Pre-attention RMSNorm} applied to the residual stream.
  \item \textbf{Grouped-query self-attention} with rotary positional embeddings (see
        Section~\ref{sec:rope}).  Projections: $W_Q \in \mathbb{R}^{d_s \times (H \cdot h_d)}$,
        $W_K \in \mathbb{R}^{d_s \times (H_k \cdot h_d)}$,
        $W_V \in \mathbb{R}^{d_s \times (H_k \cdot h_d)}$,
        $W_O \in \mathbb{R}^{(H \cdot h_d) \times d_s}$, all without bias, where
        $h_d = d_s / H$ is the per-head dimension.
  \item \textbf{Post-attention RMSNorm} added to the residual before the MLP.
  \item \textbf{MLP} (SwiGLU variant by default):
        \begin{equation}
          \mathrm{FFN}(x) = W_2\!\bigl(\sigma(W_1 x) \odot W_3 x\bigr), \quad \sigma = \mathrm{SiLU}
        \end{equation}
        where $W_1, W_3 \in \mathbb{R}^{d_{\mathrm{ff}} \times d_s}$ and
        $W_2 \in \mathbb{R}^{d_s \times d_{\mathrm{ff}}}$.
        The standard (non-gated) MLP applies $W_2(\sigma(W_1 x))$ without a gate branch.
\end{enumerate}

The backbone concludes with a final \textbf{RMSNorm} applied to the last hidden state.

\subsubsection{MLM Head}

The masked-language-modelling head converts the final hidden state into logits over the vocabulary:

\begin{equation}
  \mathrm{MLMHead}(h) = W_{\mathrm{out}}\bigl(\mathrm{LayerNorm}(\mathrm{GELU}(W_{\mathrm{in}}\, h + b_{\mathrm{in}}))\bigr) + b_{\mathrm{out}}
\end{equation}

where $W_{\mathrm{in}} \in \mathbb{R}^{d_s \times d_s}$,
$W_{\mathrm{out}} \in \mathbb{R}^{V \times d_s}$, and $V$ is the vocabulary size.

\subsubsection{Teacher Projection Layers}

When \texttt{use\_teacher\_projections = True}, the student has $N$ additional linear layers, one
per transformer layer:

\begin{equation}
  \phi^{(\ell)} : \mathbb{R}^{d_s} \to \mathbb{R}^{d_T}, \quad \phi^{(\ell)}(h) = W_{\phi}^{(\ell)} h
\end{equation}

These projections are used to match the dimensionality of the student's hidden states to the
teacher's before computing the JEPA and contrastive losses.  They are removed from the model
(and from the optimiser's parameter groups) when the EMA teacher becomes active, since the EMA
teacher already operates in student-space $\mathbb{R}^{d_s}$.

\subsection{Attention Backends}
\label{sec:attn_backends}

The attention computation is selectable via \texttt{attn\_backend}:

\begin{center}
\begin{tabular}{ll}
\toprule
\textbf{Backend} & \textbf{Implementation} \\
\midrule
\texttt{kernels\_flash} & Flash-Attention 3/2 via the \texttt{kernels} package \\
\texttt{flash\_attn}    & Flash-Attention 2 via the \texttt{flash\_attn} package \\
\texttt{flex}           & PyTorch \texttt{FlexAttention} with sliding-window + dilation mask \\
\texttt{sdpa}           & PyTorch \texttt{scaled\_dot\_product\_attention} \\
\texttt{auto}           & Resolved at runtime to the best available backend \\
\bottomrule
\end{tabular}
\end{center}

The \texttt{flex} backend applies a local sliding-window attention mask of size
\texttt{sliding\_window\_size} combined with a dilated global attention pattern with factor
\texttt{dilation}, enabling sub-quadratic attention on long sequences.

% ============================================================
\section{Tokenizer and Positional Encoding}
\label{sec:rope}
% ============================================================

\subsection{Tokenizer}

The tokenizer is loaded from \texttt{tokenizer.json} (sourced from
\texttt{Synthyra/Profluent-E1-150M}).  It is an amino-acid tokenizer with special tokens including
\texttt{<bos>}, \texttt{<eos>}, \texttt{<pad>}, and a mask token (the \texttt{?} character).
Boundary tokens \texttt{1} and \texttt{2} are used to delimit multi-chain inputs.

\subsection{Dual Rotary Positional Embedding}

Each token in the input carries two sets of position indices:

\begin{itemize}
  \item \textbf{Within-sequence position IDs} (\texttt{within\_seq\_position\_ids}): the
        position of a token within its own protein chain, up to
        \texttt{max\_num\_positions\_within\_seq} = 8\,192.
  \item \textbf{Global position IDs} (\texttt{global\_position\_ids}): the position of a token
        across the full packed context window.
\end{itemize}

Rotary embeddings are applied at each attention layer.  The within-sequence RoPE uses base
$\theta_{\mathrm{local}} = 10\,000$, while the global RoPE uses
$\theta_{\mathrm{global}} = 100\,000$.  Layers alternate between the two depending on
\texttt{global\_attention\_every\_n\_layers}: when this parameter is 0 (the default for DSM2),
all layers use within-sequence RoPE.

For an attention head dimension $h_d$, the rotary embedding for a position $p$ is the block-diagonal
matrix $R(p)$ with $2\times 2$ rotation blocks:

\begin{equation}
  R(p)_{2k,2k} = \cos\!\bigl(p\,\theta^{-2k/h_d}\bigr), \quad
  R(p)_{2k,2k+1} = -\sin\!\bigl(p\,\theta^{-2k/h_d}\bigr)
\end{equation}

for $k = 0,\ldots,h_d/2 - 1$, applied to both query and key vectors before the dot-product.

\subsection{Multi-Chain Packing}

The \texttt{E1BatchPreparer} utility handles multi-sequence inputs (e.g.\ comma-separated chains).
Up to \texttt{max\_num\_sequences} = 512 chains may be packed into one context window, each
assigned a unique sequence-type embedding from \texttt{embed\_seq\_id}
$\in \mathbb{R}^{512 \times d_s}$ that is added to the token embedding before the first layer.

% ============================================================
\section{Data Pipeline}
% ============================================================

\subsection{Dataset}

By default, the model trains on \texttt{Synthyra/uniref50}, a HuggingFace dataset providing
\texttt{train}, \texttt{valid}, and \texttt{test} splits of UniRef50 protein sequences.
The dataset is loaded via \texttt{datasets.load\_dataset}, shuffled with a fixed
\texttt{shuffle\_seed}, and optionally truncated to \texttt{train\_limit},
\texttt{valid\_limit}, and \texttt{test\_limit} samples respectively.

Each sample is a raw amino-acid string from the column specified by \texttt{sequence\_column}
(default \texttt{"sequence"}).  Sequences are tokenised and padded/truncated to at most
\texttt{max\_length} = 2\,048 tokens by the \texttt{SequenceCollator}.

\subsection{Batch Hierarchy: Patches and Patch Groups}

DSM2 uses a two-level batching scheme:

\begin{center}
\begin{tabular}{lll}
\toprule
\textbf{Level} & \textbf{Parameter} & \textbf{Default} \\
\midrule
Micro-batch (\emph{patch}) & \texttt{patch\_size} & 16 sequences \\
Patch group (per grad-accum step) & \texttt{patch\_accum} $= \lfloor\texttt{batch\_size}/\texttt{patch\_size}\rfloor$ & 4 patches \\
Gradient accumulation & \texttt{grad\_accum} & 4 steps \\
\midrule
Effective batch per optimiser step (single GPU) & \texttt{batch\_size} $\times$ \texttt{grad\_accum} & 256 sequences \\
\bottomrule
\end{tabular}
\end{center}

The DataLoader emits micro-batches of size \texttt{patch\_size}.  The trainer collects
\texttt{patch\_accum} consecutive micro-batches into a \emph{patch group}, computes the combined
loss (accumulating gradients inside \texttt{model.no\_sync()} where applicable), and repeats for
\texttt{grad\_accum} patch groups before calling \texttt{optimizer.step()}.

The effective number of sequences seen per optimiser step across all distributed ranks is:

\begin{equation}
  N_{\mathrm{eff}} = \texttt{patch\_size} \times \texttt{patch\_accum} \times \texttt{grad\_accum} \times \texttt{world\_size}
\end{equation}

\subsection{Distributed Sampling}

In multi-GPU runs (\texttt{world\_size} $> 1$), the training DataLoader uses
\texttt{DistributedSampler} with \texttt{drop\_last=True}.  Validation and test loaders also use
\texttt{DistributedSampler} with \texttt{drop\_last=False} and no shuffling.  In single-GPU runs,
\texttt{RandomSampler} is used for all splits.

On Windows, the number of workers for evaluation loaders is forced to 0 to avoid expensive
\texttt{spawn} latency; on Linux, \texttt{eval\_num\_workers} matches the training worker count.

% ============================================================
\section{Training Objective}
% ============================================================

\subsection{Stochastic Masking}
\label{sec:masking}

For each sequence $i$ in a micro-batch of size $B$, a per-sample mask rate is drawn:

\begin{equation}
  t_i \sim \mathrm{Uniform}(\varepsilon,\; 1), \qquad \varepsilon = 10^{-3}
\end{equation}

During evaluation the mask rate is fixed:

\begin{equation}
  t_i = 0.15 \quad \forall\, i
\end{equation}

A token at position $j$ (length $L$) is replaced by the mask token \texttt{[MASK]} with probability
$t_i$, subject to the constraint that special tokens (BOS, EOS, PAD, boundary tokens) and padding
positions are never masked:

\begin{equation}
  m_{ij} = \mathbf{1}\!\bigl[\mathrm{Uniform}(0,1) < t_i\bigr]
            \cdot \mathbf{1}\!\bigl[x_{ij} \notin \mathcal{S}\bigr]
            \cdot \mathrm{attn}_{ij}
\end{equation}

where $\mathcal{S}$ is the set of special token IDs and $\mathrm{attn}_{ij}$ is the attention mask.
The noisy input is $\hat{x}_{ij} = \texttt{[MASK]}$ if $m_{ij}=1$, otherwise $\hat{x}_{ij} = x_{ij}$.

The teacher always receives the \emph{original, unmasked} tokens $x_i$, while the student always
receives the masked tokens $\hat{x}_i$.

Two derived masks are computed from $m_{ij}$:

\begin{align}
  \mathcal{M} &= \{(i,j) : m_{ij} = 1\} && \text{(masked positions, CE targets)} \\
  \mathcal{D} &= \{(i,j) : m_{ij} = 0,\;\mathrm{attn}_{ij}=1,\;x_{ij}\notin\mathcal{S}\}
              && \text{(unmasked non-special positions, JEPA targets)}
\end{align}

If $\mathcal{D}$ is empty after masking (edge case), $\mathcal{D}$ falls back to all valid
attention-mask positions.  Similarly, if $\mathcal{M}$ is empty, CE falls back to all valid
attention-mask positions.

\subsection{Importance-Weighted Cross-Entropy Loss}
\label{sec:ce_loss}

The student's logits $P_\theta(x_{ij} \mid \hat{x}_i)$ are computed by the MLM head over all
positions.  Cross-entropy is evaluated only at masked positions and re-weighted by the inverse of
the sample's mask rate (importance weight $1/t_i$) to correct for the varying number of masked
tokens per sequence:

\begin{equation}
  \boxed{
    \mathcal{L}_{\mathrm{CE}} = \frac{1}{B \cdot L}
    \sum_{(i,j)\in\mathcal{M}}
    \frac{-\log P_\theta(x_{ij} \mid \hat{x}_i)}{t_i}
  }
\end{equation}

This importance weighting makes the expected gradient contribution of each sequence equal regardless
of how many tokens happen to be masked in that sample (a lower $t_i$ means fewer masked tokens but
each is up-weighted by $1/t_i$, and vice versa).

\subsection{JEPA Distillation Loss}
\label{sec:jepa_loss}

The Joint Embedding Predictive Architecture (JEPA) loss encourages the student's intermediate
representations—after projection into teacher space—to match the teacher's representations at
the unmasked context positions $\mathcal{D}$.

Let $\hat{h}_{ij}^{(\ell)} = \phi^{(\ell)}(h_{s,ij}^{(\ell)}) \in \mathbb{R}^{d_T}$ be the
$\ell$-th projected student hidden state at token $(i,j)$, and let
$h_{t,ij}^{(\ell)} \in \mathbb{R}^{d_T}$ be the corresponding teacher hidden state.

The loss applies \emph{depth weighting} $w_\ell = \ell / N$ that linearly increases with layer
depth, emphasising higher-level semantic representations:

\begin{equation}
  \boxed{
    \mathcal{L}_{\mathrm{JEPA}} =
    \frac{1}{|\mathcal{D}| \cdot N}
    \sum_{\ell=1}^{N} \frac{\ell}{N}
    \sum_{(i,j)\in\mathcal{D}}
    \frac{1}{d_T}\bigl\|\hat{h}_{ij}^{(\ell)} - h_{t,ij}^{(\ell)}\bigr\|_2^2
  }
\end{equation}

The factor $1/d_T$ comes from averaging the squared difference over the hidden dimension before
summing.  The denominator $|\mathcal{D}| \cdot N$ normalises jointly over valid token positions and
layers.  All computations are performed in \texttt{float32} for numerical stability.

\subsection{Contrastive Representation Loss}
\label{sec:contrastive_loss}

The contrastive loss aligns the \emph{relational structure} of student and teacher representations
across the batch.  Instead of matching token-level representations directly, it matches the
self-similarity matrices of pooled sequence representations at each layer.

\textbf{Step 1: Mean-variance pooling.}
For each layer $\ell$ and sample $i$, the student and teacher hidden states are pooled over valid
(attended) token positions into a $2d_T$-dimensional vector:

\begin{equation}
  \bar{s}_i^{(\ell)} = \bigl[\mu\bigl(\hat{H}_i^{(\ell)}\bigr);\; \sigma^2\bigl(\hat{H}_i^{(\ell)}\bigr)\bigr]
  \in \mathbb{R}^{2d_T}
\end{equation}

where $\mu$ and $\sigma^2$ are the masked mean and variance over the sequence dimension.
This yields tensors $\bar{s}^{(\ell)}, \bar{t}^{(\ell)} \in \mathbb{R}^{B \times 2d_T}$.

\textbf{Step 2: Self-similarity matrices.}
For each layer, compute the normalised intra-batch similarity matrices:

\begin{equation}
  S_{ij}^{(\ell)} = \mathrm{softmax}\!\bigl(\bar{s}^{(\ell)}\,\bar{s}^{(\ell)\top}\bigr)_{ij},
  \qquad
  T_{ij}^{(\ell)} = \mathrm{softmax}\!\bigl(\bar{t}^{(\ell)}\,\bar{t}^{(\ell)\top}\bigr)_{ij}
\end{equation}

where softmax is applied row-wise over the $B$ columns, yielding
$S^{(\ell)}, T^{(\ell)} \in \mathbb{R}^{B \times B}$.

\textbf{Step 3: Depth-weighted MSE.}
The loss is the depth-weighted mean squared error between the two similarity matrices, averaged
over layers and batch pairs:

\begin{equation}
  \boxed{
    \mathcal{L}_{\mathrm{contrastive}} =
    \frac{1}{B^2 \cdot N}
    \sum_{\ell=1}^{N} \frac{\ell}{N}
    \sum_{i=1}^{B} \sum_{j=1}^{B}
    \bigl(S_{ij}^{(\ell)} - T_{ij}^{(\ell)}\bigr)^2
  }
\end{equation}

Note: the contrastive loss operates over the \emph{full patch group} (all patches concatenated),
not per micro-batch, so $B$ here is the effective group batch size.  All computations are in
\texttt{float32}.

\subsection{Auxiliary Loss Warmup Scale}
\label{sec:aux_warmup}

The auxiliary losses (JEPA and contrastive) are gated by a scalar warmup factor $\gamma \in [0,1]$
that depends on the current optimiser step $s$:

\begin{equation}
  \gamma(s) =
  \begin{cases}
    0 & s < T_{\mathrm{free}} \\[4pt]
    \dfrac{s - T_{\mathrm{free}}}{T_{\mathrm{warm}}} & T_{\mathrm{free}} \le s < T_{\mathrm{free}} + T_{\mathrm{warm}} \\[8pt]
    1 & s \ge T_{\mathrm{free}} + T_{\mathrm{warm}}
  \end{cases}
\end{equation}

where:
\begin{align*}
  T_{\mathrm{free}} &= \lfloor S_{\max} \times \texttt{teacher\_free\_percent} \rfloor \\
  T_{\mathrm{warm}} &= \lfloor S_{\max} \times \texttt{aux\_loss\_warmup\_percent} \rfloor
\end{align*}

The effective auxiliary weights are $\tilde\alpha_{\mathrm{JEPA}} = \gamma \cdot \alpha_{\mathrm{JEPA}}$
and $\tilde\alpha_{\mathrm{contrastive}} = \gamma \cdot \alpha_{\mathrm{contrastive}}$.

\subsection{Auxiliary-to-CE Ratio Guard}

To prevent instabilities caused by auxiliary losses dominating the CE signal, an adaptive scaling
factor $\kappa \in (0, 1]$ is applied to the total weighted auxiliary loss whenever it exceeds a
specified fraction $r_{\max}$ of the CE loss magnitude:

\begin{equation}
  A = \bigl|\tilde\alpha_{\mathrm{JEPA}}\,\mathcal{L}_{\mathrm{JEPA}}
          + \tilde\alpha_{\mathrm{contrastive}}\,\mathcal{L}_{\mathrm{contrastive}}\bigr|, \qquad
  C = \bigl|\alpha_{\mathrm{CE}}\,\mathcal{L}_{\mathrm{CE}}\bigr|
\end{equation}

\begin{equation}
  \kappa =
  \begin{cases}
    r_{\max} \cdot C \,/\, A & A > r_{\max} \cdot C \text{ and } A > 0 \\
    1 & \text{otherwise}
  \end{cases}
\end{equation}

Default $r_{\max} = 0.25$ (set \texttt{max\_aux\_to\_ce\_ratio} $\le 0$ to disable the guard).

\subsection{Combined Loss Function}

The total scalar loss for a patch group is:

\begin{equation}
  \boxed{
    \mathcal{L} =
    \alpha_{\mathrm{CE}}\,\mathcal{L}_{\mathrm{CE}}
    + \kappa\!\left(\tilde\alpha_{\mathrm{JEPA}}\,\mathcal{L}_{\mathrm{JEPA}}
                  + \tilde\alpha_{\mathrm{contrastive}}\,\mathcal{L}_{\mathrm{contrastive}}\right)
  }
\end{equation}

Each loss component is accumulated in a weighted sum over the patch group; the weight of patch $p$
with $B_p$ sequences is $B_p / B_{\mathrm{total}}$ where $B_{\mathrm{total}} = \sum_p B_p$.
CE and JEPA contributions are accumulated per-patch; the contrastive loss is computed once over the
full concatenated batch of pooled representations.

% ============================================================
\section{Training Schedule}
% ============================================================

\subsection{Three-Phase Loss Schedule}

Training proceeds through three consecutive phases controlled by \texttt{teacher\_free\_percent}
and \texttt{aux\_loss\_warmup\_percent}:

\begin{enumerate}
  \item \textbf{Teacher-free phase} (steps $[0, T_{\mathrm{free}})$, default 10\% of total): only
        the CE loss is active ($\gamma = 0$).  The teacher is not forwarded, saving compute.

  \item \textbf{Auxiliary warmup phase} (steps $[T_{\mathrm{free}}, T_{\mathrm{free}}+T_{\mathrm{warm}})$,
        default next 10\%): $\gamma$ ramps linearly from 0 to 1.  The teacher is forwarded and
        JEPA + contrastive losses are gradually introduced.

  \item \textbf{Full training phase} (steps $\ge T_{\mathrm{free}}+T_{\mathrm{warm}}$, default
        remaining 80\%): all three losses are active at full weight, subject only to the ratio guard.
\end{enumerate}

\subsection{EMA Teacher Transition}
\label{sec:ema}

At step $\lfloor S_{\max} \times \texttt{ema\_start\_percent}\rfloor$ (default 40\% of total
steps), the \texttt{EMATeacherCallback} activates.  A deep copy of the current student weights is
made, frozen, compiled with \texttt{torch.compile}, and attached to the model as
\texttt{model.ema\_teacher}.

Thereafter, at the end of every optimiser step, the EMA teacher parameters $\theta_T$ are updated:

\begin{equation}
  \theta_T \leftarrow \delta\,\theta_T + (1 - \delta)\,\theta_S
\end{equation}

where $\delta = \texttt{ema\_decay}$ (default 0.999) and $\theta_S$ are the current student
parameters.

Once the EMA teacher is active, the trainer switches the active teacher from the external frozen
model to the EMA copy, and performs a one-time cleanup:

\begin{itemize}
  \item The teacher projection layers $\{\phi^{(\ell)}\}$ are removed from both the student and the
        EMA teacher (they are no longer needed since the EMA teacher lives in student-space $\mathbb{R}^{d_s}$).
  \item The corresponding parameters are removed from the optimiser's parameter groups.
  \item The original external teacher model is set to \texttt{None} and its GPU memory is freed via
        \texttt{torch.cuda.empty\_cache()}.
\end{itemize}

The EMA teacher is forwarded by calling the student's own \texttt{forward} with
$\alpha_{\mathrm{CE}} = \alpha_{\mathrm{JEPA}} = \alpha_{\mathrm{contrastive}} = 0$, retrieving
\texttt{student\_hidden\_states} from the output rather than hidden states from the transformer backbone.

\subsection{Learning Rate Schedule}

The learning rate follows a three-phase piecewise schedule defined as a multiplier on the base
learning rate \texttt{lr}:

\begin{equation}
  \lambda(s) =
  \begin{cases}
    \dfrac{s+1}{W}
      & 0 \le s < W \quad (\text{linear warmup},\ W = \lceil 0.01\,S_{\max} \rceil) \\[8pt]
    1.0
      & W \le s < W + P \quad (\text{plateau},\ P = \lceil 0.49\,S_{\max} \rceil) \\[8pt]
    \dfrac{1}{2}\!\left(1 + \cos\!\left(\pi \cdot \dfrac{s - (W+P) + 1}{C}\right)\right)
      & W+P \le s < W+P+C \quad (\text{cosine cooldown},\ C = \lceil 0.30\,S_{\max} \rceil) \\[8pt]
    0
      & s \ge W + P + C
  \end{cases}
\end{equation}

The schedule is implemented as \texttt{torch.optim.lr\_scheduler.LambdaLR}.  The remaining $\approx
20\%$ of steps after the cooldown have zero learning rate.

% ============================================================
\section{Optimisation}
% ============================================================

\subsection{Standard AdamW Path}

When \texttt{--muon} is not supplied, all student parameters are optimised with AdamW:

\begin{equation}
  \texttt{AdamW}(\theta,\; \text{lr}=\texttt{lr},\; \text{weight\_decay}=0)
\end{equation}

\subsection{MuonClip Path}

When \texttt{--muon} is supplied, parameters are partitioned into three groups:

\begin{center}
\begin{tabular}{lll}
\toprule
\textbf{Group} & \textbf{Selection rule} & \textbf{Optimiser} \\
\midrule
\texttt{muon\_params} & \texttt{ndim} $\ge 2$, name does not contain \texttt{embed},
                        \texttt{lm\_head}, or \texttt{mlm\_head} & MuonClip (Muon + QK-Clip) \\
\texttt{adamw\_params} & embeddings and LM-head weight matrices & AdamW \\
\texttt{attention\_params} & $(W_Q, W_K)$ pairs from every attention layer (QK-Clip targets) & (used inside MuonClip) \\
\bottomrule
\end{tabular}
\end{center}

The \texttt{MuonAdamWWrapper} combines both sub-optimisers behind a single
\texttt{torch.optim.Optimizer} interface.

\subsubsection{Muon (Momentum + Orthogonalisation)}

The Muon update for a weight matrix $W \in \mathbb{R}^{m \times n}$ applies Nesterov momentum and
then orthogonalises the update via five steps of the Newton–Schulz iteration before scaling by the
learning rate:

\begin{equation}
  \tilde{G} = \mathrm{NewtonSchulz5}(G_t), \qquad W \leftarrow W - \eta_{\mathrm{muon}}\,\tilde{G}
\end{equation}

where $G_t$ is the current (momentum-adjusted) gradient and $\mathrm{NewtonSchulz5}$ maps $G_t$ to
a matrix with approximately orthonormal rows (or columns), stabilising the effective learning rate
across layers of different widths.

\subsubsection{QK-Clip}

For each attention layer $\ell$ and head $h$, the maximum attention logit magnitude observed during
the most recent forward pass is tracked:

\begin{equation}
  s_{\max}^{(\ell, h)} = \max_{i,j} \frac{\langle q_i^{(\ell,h)},\, k_j^{(\ell,h)} \rangle}{\sqrt{h_d}}
\end{equation}

where the max is over sequence positions.  When \texttt{output\_s\_max = True} (default during
training), the model collects these values and passes them back as \texttt{DSM2Output.s\_max}.  The
trainer aggregates $s_{\max}$ across patches by taking the element-wise maximum:

\begin{equation}
  s_{\max}^{(\ell,h)} \leftarrow \max_{p}\; s_{\max,p}^{(\ell,h)}
\end{equation}

Before the optimiser step, $s_{\max}$ is injected into the \texttt{MuonClip} instance.  QK-Clip
then limits the effective step size on $W_Q$ and $W_K$ to prevent the attention logit magnitude from
growing beyond the threshold $\tau = \texttt{muon\_tau}$ (default 100).

\subsubsection{Gradient Clipping}

When \texttt{max\_grad\_norm} $> 0$, global gradient clipping is applied to the student model
before \texttt{optimizer.step()}:

\begin{equation}
  g \leftarrow g \cdot \min\!\left(1,\; \frac{\texttt{max\_grad\_norm}}{\|g\|_2}\right)
\end{equation}

% ============================================================
\section{Distributed Training}
% ============================================================

\subsection{Process Group Initialisation}

If \texttt{world\_size} $> 1$ and \texttt{--no\_init\_distributed} is not set, the trainer calls
\texttt{torch.distributed.init\_process\_group} with the specified \texttt{distributed\_backend}
(default \texttt{gloo}) and \texttt{init\_method="env://"}.  Environment variables
\texttt{RANK}, \texttt{WORLD\_SIZE}, and \texttt{LOCAL\_RANK} are read by
\texttt{infer\_rank\_world\_size\_local\_rank}.

\subsection{Model Parallelism}

After \texttt{torch.compile}, the model is wrapped in
\texttt{DistributedDataParallel} (DDP) on CUDA runs:

\begin{verbatim}
DistributedDataParallel(model, device_ids=[local_rank], output_device=local_rank)
\end{verbatim}

Gradient synchronisation across ranks is deferred during inner patch-group accumulation steps using
\texttt{model.no\_sync()}.  A full \texttt{all\_reduce} occurs only at the final patch of each
gradient accumulation step.

\subsection{Distributed Loss Reduction}

During evaluation, per-rank losses are reduced via \texttt{dist.all\_reduce} (sum, then divide by
world size).  Logits and labels are gathered across ranks with
\texttt{dist.all\_gather\_object} and concatenated on the main process (rank 0) before metric
computation.

% ============================================================
\section{Evaluation}
% ============================================================

\subsection{Evaluation Protocol}

An initial evaluation on the test set is performed before training begins.  During training,
evaluation runs on the validation set every \texttt{eval\_every} optimiser steps (defaulting to
\texttt{save\_every} = 1\,000 when not specified).  A final evaluation on the test set is performed
after training completes.

\subsection{Metrics}

The \texttt{ComputeDSM2Metrics} class computes the following metrics from the predicted logits
$\hat{Y} \in \mathbb{R}^{B \times L \times V}$ and ground-truth token IDs $Y \in \mathbb{Z}^{B
\times L}$ (positions with label $-100$ are ignored):

\begin{center}
\begin{tabular}{ll}
\toprule
\textbf{Metric} & \textbf{Definition / Notes} \\
\midrule
\textbf{Cross-entropy loss} & $\mathcal{L}_{\mathrm{CE}}$ computed in \texttt{float32} over masked positions \\
\textbf{Alignment score}   & Per-sequence alignment score from \texttt{GetAlignmentScoreFromLogits}, averaged over batch \\
\textbf{Accuracy}          & Top-1 accuracy $\hat{y}_{ij} = \arg\max_v \hat{Y}_{ijv}$ at masked positions \\
\textbf{F1 (weighted)}     & Weighted-average F1 over residue classes at masked positions \\
\textbf{Precision (weighted)} & Weighted-average precision at masked positions \\
\textbf{Recall (weighted)} & Weighted-average recall at masked positions \\
\textbf{MCC}               & Matthews correlation coefficient at masked positions \\
\bottomrule
\end{tabular}
\end{center}

All classification metrics are computed with \texttt{scikit-learn} functions after flattening
valid (non-$-100$) positions.

% ============================================================
\section{Checkpointing and Logging}
% ============================================================

\subsection{Checkpointing}

Every \texttt{save\_every} optimiser steps, the main process (rank 0) saves:

\begin{itemize}
  \item The unwrapped student model via \texttt{save\_pretrained} into
        \texttt{<save\_path>/checkpoint-step-<N>/}.
  \item Optimiser state to \texttt{optimizer.pt}.
  \item Scheduler state to \texttt{scheduler.pt}.
\end{itemize}

At the end of training, when an HF token is provided, the final model is pushed to the HuggingFace
Hub at the path specified by \texttt{save\_path} (private repository).

\subsection{Weights \& Biases Logging}

When both \texttt{WANDB\_AVAILABLE=true} and \texttt{--wandb\_token} are provided, the following
are logged every \texttt{logging\_steps} optimiser steps under the \texttt{train/} prefix and after
each evaluation under \texttt{valid/} or \texttt{test/} prefixes:

\begin{center}
\begin{tabular}{ll}
\toprule
\textbf{Key} & \textbf{Description} \\
\midrule
\texttt{ce\_loss} & Raw (unweighted) CE loss for the patch group \\
\texttt{jepa\_loss} & Raw JEPA loss \\
\texttt{contrastive\_loss} & Raw contrastive loss \\
\texttt{weighted\_ce\_loss} & $\alpha_{\mathrm{CE}} \cdot \mathcal{L}_{\mathrm{CE}}$ \\
\texttt{weighted\_jepa\_loss} & $\tilde\alpha_{\mathrm{JEPA}} \cdot \mathcal{L}_{\mathrm{JEPA}}$ \\
\texttt{weighted\_contrastive\_loss} & $\tilde\alpha_{\mathrm{contrastive}} \cdot \mathcal{L}_{\mathrm{contrastive}}$ \\
\texttt{weighted\_aux\_loss\_pre\_guard} & Weighted aux sum before ratio guard \\
\texttt{weighted\_aux\_loss} & Weighted aux sum after ratio guard \\
\texttt{aux\_loss\_scale} & Current $\gamma(s)$ \\
\texttt{aux\_guard\_scale} & Current $\kappa$ \\
\bottomrule
\end{tabular}
\end{center}

% ============================================================
\section{Hyperparameters}
% ============================================================

Table~\ref{tab:hparams} lists every command-line argument, its type, default value, and purpose.

\begin{longtable}{p{5.2cm} p{1.3cm} p{3.2cm} p{5.0cm}}
\caption{All DSM2 training arguments and their defaults.\label{tab:hparams}} \\
\toprule
\textbf{Argument} & \textbf{Type} & \textbf{Default} & \textbf{Description} \\
\midrule
\endfirsthead
\multicolumn{4}{c}{\tablename~\thetable{} (continued)} \\
\toprule
\textbf{Argument} & \textbf{Type} & \textbf{Default} & \textbf{Description} \\
\midrule
\endhead
\bottomrule
\endfoot
\bottomrule
\endlastfoot

% Runtime
\texttt{--hf\_token} & str & \texttt{None} & HuggingFace authentication token \\
\texttt{--wandb\_token} & str & \texttt{None} & Weights \& Biases API token \\
\texttt{--wandb\_project} & str & \texttt{DSM2} & W\&B project name \\
\texttt{--save\_path} & str & \texttt{GleghornLab/DSM2\_600} & HF Hub repo and local checkpoint prefix \\
\texttt{--distributed\_backend} & str & \texttt{gloo} & \texttt{torch.distributed} backend \\
\texttt{--no\_init\_distributed} & flag & \texttt{False} & Skip \texttt{init\_process\_group} \\
\texttt{--no\_pin\_memory} & flag & \texttt{False} & Disable pinned-memory DataLoaders \\
\texttt{--bugfix} & flag & \texttt{False} & Use tiny debug configuration \\
\midrule

% Model
\texttt{--teacher\_model\_path} & str & \texttt{Synthyra/DPLM-3B} & HF path for the teacher model \\
\texttt{--pretrained\_weights} & str & \texttt{None} & HF path to seed student weights; skips projection layers if set \\
\texttt{--student\_hidden\_size} & int & \texttt{256} & $d_s$, must satisfy $d_s \bmod 64 = 0$ \\
\texttt{--student\_expansion\_ratio} & float & \texttt{4.0} & $d_{\mathrm{ff}} / d_s$ ratio for the MLP \\
\texttt{--sliding\_window\_size} & int & \texttt{512} & Sliding-window size for \texttt{flex} attention \\
\texttt{--dilation} & int & \texttt{16} & Dilation factor for \texttt{flex} attention \\
\texttt{--attn\_backend} & str & \texttt{auto} & Attention backend (see Section~\ref{sec:attn_backends}) \\
\midrule

% Data
\texttt{--data\_path} & str & \texttt{Synthyra/uniref50} & HuggingFace dataset identifier \\
\texttt{--sequence\_column} & str & \texttt{sequence} & Column name containing amino-acid sequences \\
\texttt{--max\_length} & int & \texttt{2048} & Maximum tokenised sequence length \\
\texttt{--train\_limit} & int & \texttt{0} & Max train samples ($\le 0$ = all) \\
\texttt{--valid\_limit} & int & \texttt{1000} & Max validation samples \\
\texttt{--test\_limit} & int & \texttt{1000} & Max test samples \\
\texttt{--shuffle\_seed} & int & \texttt{42} & Random seed for dataset shuffling \\
\texttt{--dataloader\_num\_workers} & int & \texttt{0} & DataLoader worker processes \\
\texttt{--dataloader\_prefetch\_factor} & int & \texttt{2} & Prefetch factor when workers $> 0$ \\
\midrule

% Optimisation
\texttt{--lr} & float & \texttt{1e-4} & Base learning rate for AdamW \\
\texttt{--batch\_size} & int & \texttt{64} & Effective batch size metadata; sets \texttt{patch\_accum} $= \lfloor\texttt{batch\_size}/\texttt{patch\_size}\rfloor$ \\
\texttt{--patch\_size} & int & \texttt{16} & Micro-batch size per forward pass \\
\texttt{--grad\_accum} & int & \texttt{4} & Gradient accumulation steps \\
\texttt{--max\_steps} & int & \texttt{100000} & Total optimiser steps $S_{\max}$ \\
\texttt{--max\_grad\_norm} & float & \texttt{0.0} & Gradient clipping norm ($0$ = disabled) \\
\texttt{--muon} & flag & \texttt{False} & Enable MuonClip optimiser \\
\texttt{--muon\_lr} & float & \texttt{0.001} & Learning rate for the Muon sub-optimiser \\
\texttt{--muon\_tau} & float & \texttt{100.0} & QK-Clip $\tau$ threshold \\
\midrule

% Loss
\texttt{--alpha\_ce} & float & \texttt{1.0} & Weight $\alpha_{\mathrm{CE}}$ on CE loss \\
\texttt{--alpha\_jepa} & float & \texttt{0.1} & Weight $\alpha_{\mathrm{JEPA}}$ on JEPA loss \\
\texttt{--alpha\_contrastive} & float & \texttt{1.0} & Weight $\alpha_{\mathrm{contrastive}}$ on contrastive loss \\
\texttt{--teacher\_free\_percent} & float & \texttt{0.1} & Fraction of steps in CE-only phase \\
\texttt{--aux\_loss\_warmup\_percent} & float & \texttt{0.1} & Fraction of steps for aux-loss linear warmup \\
\texttt{--max\_aux\_to\_ce\_ratio} & float & \texttt{0.25} & Cap $r_{\max}$ on aux/CE ratio ($\le 0$ = disabled) \\
\midrule

% EMA
\texttt{--ema\_start\_percent} & float & \texttt{0.4} & Fraction of steps before EMA teacher activates \\
\texttt{--ema\_decay} & float & \texttt{0.999} & EMA decay $\delta$ \\
\midrule

% Logging / checkpointing
\texttt{--save\_every} & int & \texttt{1000} & Save checkpoint every $N$ optimiser steps \\
\texttt{--eval\_every} & int & \texttt{0} & Evaluate every $N$ steps ($\le 0$ uses \texttt{save\_every}) \\
\texttt{--logging\_steps} & int & \texttt{100} & Log training metrics every $N$ steps \\
\end{longtable}

% ============================================================
\section{Summary of Training Phases}
% ============================================================

Figure~\ref{fig:phases} summarises the interplay of all training phases across the total number of
optimiser steps $S_{\max}$ under the default hyperparameter settings.

\begin{center}
\small
\begin{tabular}{ccccc}
\toprule
\textbf{Step range} & \textbf{\% of $S_{\max}$} & \textbf{CE} & \textbf{JEPA + Contrastive} & \textbf{Teacher source} \\
\midrule
$[0,\; 0.10\,S_{\max})$   & 10\% & \checkmark & $\times$ & (not forwarded) \\
$[0.10,\; 0.20\,S_{\max})$ & 10\% & \checkmark & linear ramp $\gamma \in [0,1]$ & External (frozen) \\
$[0.20,\; 0.40\,S_{\max})$ & 20\% & \checkmark & \checkmark (full) & External (frozen) \\
$[0.40,\; 1.00\,S_{\max})$ & 60\% & \checkmark & \checkmark (full) & EMA (student self-distillation) \\
\bottomrule
\end{tabular}
\end{center}

\captionof*{table}{Figure~1: Training phase schedule under default settings.  At step 0.40$\,S_{\max}$ the
EMA teacher activates and the external teacher + projection layers are discarded.}
\label{fig:phases}

\end{document}
